Predicting the direction of a target gene after perturbation does not establish that a model can identify the underlying regulatory mechanism. We introduce a benchmark for transcription factor (TF)-to-gene regulation that integrates mechanism-specific experimental evidence. We evaluate whether models make accurate predictions and provide concise mechanistic explanations. We developed a benchmark comprising three tasks. Task 1 assesses mechanism classification using 171 TF-to-gene pairs assigned to four regulatory categories: direct cis-regulatory DNA, chromatin-mediated regulation, combinatorial TF regulation and regulatory RNA-mediated regulation. We curated regulatory mechanisms and then matched perturbations for task 2, yielding 193 forward-reasoning units, and constructed task 3 that includes inverse-intervention units from different perturbation types for the same TF-to-gene pairs. Tasks 2 and 3 also evaluate the mechanistic explanations supporting model predictions. Explanations are assessed for mechanism, mediator, causal role and order, intervention-to-outcome logic, and consistency with curated evidence. Across six evaluated systems, Task 1 macro-F1 and Task 2/3 joint F1 scores range from 22.7\% to 34.5\% for Task 1, 57.5\% to 78.4\% for Task 2, and 25.6\% to 35.7\% for Task 3. Full-chain reasoning pass rates range from 62.0\% to 84.7\% in forward reasoning of Task 2 and 43.8\% to 72.5\% in inverse reasoning of Task 3. The evaluated systems show recurring difficulties in recovering required regulatory dependencies and their context-specific causal roles. These results identify specific weaknesses in evidence-aligned mechanistic explanations under the benchmark conditions. 
%Our benchmark provides a focused behavioral assessment of biological mechanism understanding.

\par\medskip
\noindent\textbf{Keywords:} TF-to-gene regulation; perturbation responses; LLMs' mechanistic reasoning
